%%%%%%%%%%%%%%%%%%%%%%%%%%%%%%%%%%%%%%%%%%%%%%%%%%%%%%%%%%%%%%%%%%%%%%
%  第 4 个逻辑月 · 月度报告
%  主题：网络 · 优化 · 连接 · UDF 修正
%
%  说明：本月大部分工作尚未开展，报告以「已知条件 + 待办清单」形式组织。
%  编译：xelatex 本文件（跑两遍）
%%%%%%%%%%%%%%%%%%%%%%%%%%%%%%%%%%%%%%%%%%%%%%%%%%%%%%%%%%%%%%%%%%%%%%
\documentclass[stage2,fontset=fandol]{ydzjreport}

\stagetwosetup
  {第 4 个逻辑月 · 连接成网络}%
  {冯旭东、陈思文、张昊洋}%
  {复合微通道分区冷却}%
  {【待确认】逻辑月 4 对应的自然月}%
  {网络拓扑；连接段；分流优化；UDF 修正收尾}%
  {把单元级成果推进到整板流路网络：设计小块连接处的过渡与分流通道，
   并完成新物性 UDF 的挂接与校验。}%

\begin{document}
\makestagetwohead

\begin{notebox}
\textbf{本月完成情况说明}：截至整理本报告时，本逻辑月的四项工作
（网络、优化、连接、UDF 修正）\textbf{仅"连接"有方向性方案，
其余尚未开展}。本报告如实反映这一状态，未开展的部分按规范留空标注。
\end{notebox}

% ==================================================================
\section{本月研究任务}

\subsection{任务目标}

\begin{enumerate}
  \item \textbf{网络}：确定小块之间连成整板的流路拓扑（串联、并联或混合）；
  \item \textbf{连接}：设计小块连接处的过渡与分流通道；
  \item \textbf{优化}：按各区冷却需求合理分配流量；
  \item \textbf{UDF 修正}：完成新物性 UDF 的挂接、物性自检与代表工况复算。
\end{enumerate}

\subsection{完成情况}

\begin{table}[H]
\centering
\caption{本月任务完成情况}
\label{tab:tasks}
\begin{tabularx}{\textwidth}{@{}Y c Y@{}}
\toprule
\rowcolor{ydzjlight} 任务 & 状态 & 说明 \\
\midrule
流路网络拓扑 & 未开始 & 【待补：网络拓扑方案】 \\
小块连接段设计 & 部分完成 & 已明确设计方向与评价方法，方案待出 \\
支路流量分配优化 & 未开始 & 【待补：分流方案与流量数据】 \\
UDF 修正收尾 & 部分完成 & 编译已完成，挂接与复算待推进 \\
\bottomrule
\end{tabularx}
\end{table}

【待确认5】本月「连接」与上逻辑月「形态」的边界。
【待确认6】UDF 编译已提前完成，本月该项是否仍按原计划执行。

% ==================================================================
\section{进阶阅读与理论支撑}

\subsection{本阶段文献与资料}

\begin{table}[H]
\centering
\caption{本阶段文献与资料}
\label{tab:lit}
\begin{tabularx}{\textwidth}{@{}c Y Y@{}}
\toprule
\rowcolor{ydzjlight} 序号 & 文献与资料 & 解决了我哪个问题 \\
\midrule
1 & 申报书中复合冷板的分区方案 & 确定网络拓扑的总体约束 \\
2 & 新水物性 UDF 源码 & 明确挂接与校验所需的条件 \\
3 & 流动分配不均匀性的相关研究 & 理解并联支路供液不足的机理 \\
\bottomrule
\end{tabularx}
\end{table}

\subsection{对本方向的启示}

\textbf{一、单元性能不等于整板性能。}
这是本月最重要的认识。单元级的仿真结果是在
"独立、充分供液"的假设下得到的；一旦连成整板：

\begin{itemize}
  \item \textbf{并联}时，高阻力支路可能\textbf{供液不足}，
        原本温度最低的流道反而得不到足够流量；
  \item \textbf{串联}时，上游吸热会\textbf{抬高下游入口水温}，
        使下游小块的冷却能力下降。
\end{itemize}

因此上逻辑月得到的单元排序\textbf{不能直接搬到整板}。

\textbf{二、连接段是新的阻力源。}
小块之间的过渡与分流通道会引入额外压降，
且容易在局部产生回流与偏流。连接段设计的评价，应同时看
\textbf{入口速度分布的均匀性}与\textbf{连接段压降}。

% ==================================================================
\section{研究进展}

\subsection{已完成的工作}

\textbf{（一）明确了连接段的设计方向与评价方法}

已确定的\textbf{设计目标}：优化连接段截面及分流布局，
减少局部回流、流动突变和偏流，使各小块入口的速度分布尽量均匀、
流动整体保持稳定。

已确定的\textbf{评价方法}：结合入口速度分布、支路流量及连接段压降评价效果，
并按各区冷却需求合理分配流量。

\textbf{（二）UDF 修正的进度}

上逻辑月已完成新水物性 UDF 的编写与编译（含温度范围检查、
物性自检输出、比热与焓的一致性处理）。本月的工作是\textbf{挂接与校验}。

\subsection{关键结果}

【待补：本月尚无仿真结果】

本月四项工作中，连接段方案、网络拓扑、分流优化均未形成仿真结果，
故此处无量化数据可填。已具备的\textbf{前置条件}如下：

\begin{table}[H]
\centering
\caption{本月工作的前置条件}
\label{tab:ready}
\begin{tabularx}{\textwidth}{@{}Y Y Y@{}}
\toprule
\rowcolor{ydzjlight} 前置条件 & 状态 & 来源 \\
\midrule
单元级网格基准 & 已有 & 上逻辑月网格无关性检验，约 150 万 \\
单元级性能数据 & 已有 & 四种流道 39 个工况 \\
新物性 UDF & 已编译 & 上逻辑月完成 \\
连接段设计方向 & 已明确 & 本月 \\
整板分区方案 & 【待确认】 & 取决于小块划分定义 \\
\bottomrule
\end{tabularx}
\end{table}

\subsection{与预期的差异}

【待补：本月无仿真结果，暂无可报告的差异】

\textbf{可以报告的是一项认识上的变化}：
原以为单元级排序可以直接用于整板设计，
现在认识到\textbf{必须经过"连接与分流"这一环节的检验}——
单元最优不等于整板最优。

% ==================================================================
\section{问题与对策}

\subsection{遇到的主要问题}

\begin{enumerate}
  \item \textbf{整板分区方案未定。} 整板的网络拓扑取决于"小块如何划分"，
        而上逻辑月的划分定义【待确认9】至今未统一，
        导致网络设计缺少输入。
  \item \textbf{连接段的优化目标尚未量化。} 目前已明确"减少回流与偏流"的方向，
        但"入口速度分布多均匀算合格"没有判据。
  \item \textbf{UDF 挂接的具体条件未确认。} 需要确认入口及实际流体温区是否
        落在源码的 280 至 360 K 检查范围内。
\end{enumerate}

\subsection{已尝试的方案}

\begin{enumerate}
  \item 针对问题 1，已梳理现有分区说法（低温区／变化区／高温区），
        但未能统一到一份定义。
  \item 针对问题 2，拟采用"入口速度分布的均匀性指标 + 连接段压降"作为双重判据，
        具体阈值待与前三项工作对齐后确定。
  \item 针对问题 3，已确认新 UDF 源码自带温度范围检查，
        挂接后可直接读取自检输出。
\end{enumerate}

\subsection{下一步对策}

\begin{enumerate}
  \item 先确定小块划分定义，再据此设计网络拓扑；
  \item 连接段设计同时输出\textbf{阻力系数或压降—流量关系}，
        既服务于分流优化，也可作为控制方向"执行器到流速"函数的管路阻抗输入；
  \item 挂接新 UDF，核查入口及实际流体温区，完成物性自检与代表工况复算。
\end{enumerate}

\begin{keybox}
\textbf{与控制方向的接口}：连接段的阻力特性即控制模型中
"执行器到流速"函数里的\textbf{管路阻抗}。
建议连接段设计\textbf{同时产出阻力系数}，一份工作两用。
\end{keybox}

% ==================================================================
\section{阶段性结论}

\subsection{可以确定的}

\begin{enumerate}
  \item 已明确连接段的设计目标与评价方法（入口速度分布均匀性 + 连接段压降）。
  \item 已具备开展网络设计的前置条件：网格基准、单元性能数据、新物性 UDF。
  \item \textbf{单元性能不能直接代表整板性能}——
        并联供液不足与串联入口水温抬升是两个必须处理的效应。
\end{enumerate}

\subsection{尚不能确定的}

\begin{enumerate}
  \item \textbf{网络拓扑的形式}（串联／并联／混合），
        缺少的证据是：小块划分定义与分区方案。
  \item \textbf{分流是否合理}，
        缺少的证据是：各支路流量与连接段压降的仿真数据。
  \item \textbf{物性修正对整板结果的影响}，
        缺少的证据是：挂接新 UDF 后的复算结果。
\end{enumerate}

% ==================================================================
\section{下月计划}

\begin{table}[H]
\centering
\caption{下月工作计划}
\label{tab:plan}
\begin{tabularx}{\textwidth}{@{}Y Y c@{}}
\toprule
\rowcolor{ydzjlight} 事项 & 预期产出 & 时间 \\
\midrule
确定小块划分定义，据此设计网络拓扑 & 网络拓扑方案 & 第 1 周 \\
连接段过渡与分流通道建模 & 连接段形态方案与入口流场对比 & 第 2--3 周 \\
支路流量分配与连接段压降校核 & 分流方案与压降数据 & 第 3 周 \\
整板组合模型与初步分流核对 & 分区布局与调控基础数据 & 第 4 周 \\
\bottomrule
\end{tabularx}
\end{table}

% ==================================================================
\begin{thebibliography}{9}
\bibitem{proposal} 大创项目申报书：研究内容与实施方案. 2026.
\bibitem{udf} 项目组. water\_properties\_refprop\_fit.c：液态水物性拟合与 Fluent UDF 源码. 项目计算资料, 2026.
\bibitem{data} 四种流道数据与折线图. 项目计算资料, 2026.
\end{thebibliography}

\end{document}
